\section[Fiber-categories]{Fiber-categories; equivalence of
$\cal{E}$-categories}
\label{VI.4}

Let~$\cal{F}$ be a category over~$\cal{E}$, and let $S \in
\Ob({\cal{E}})$. One calls \emph{fiber-category of}~$\cal{F}$
at~$S$
the subcategory~$\cal{F}_S$
of~$\cal{F}$, inverse image of the
% original p. 159
punctual subcategory of~$\cal{E}$ defined by $S$. Thus
the objects of~$\cal{F}_S$ are the objects~$\xi$ of~$\cal{F}$ such that
$p(\xi)=S$, its morphisms are the morphisms~$u$ of~$\cal{F}$ such
that $p(u)=\id_S$, \ie the $S$-morphisms in $\cal{F}$. Of
course,~$\cal{F}_S$ is canonically isomorphic to the fibered product
$\cal{F} \times_{\cal{E}} \{S\}$, where~$\{S\}$ denotes the
punctual subcategory of~$\cal{E}$ defined by~$S$, equipped
with its inclusion functor into~$\cal{E}$. It follows (taking
into account the transitivity of change of base) that if one performs the
change of base $\lambda\colon \cal{E}' \to \cal{E}$, then for
every $S' \in \Ob(\cal{E}')$, \emph{the projection $\pr_1\colon
\cal{F}' = \cal{F} \times_{\cal{E}} \cal{E}' \to \cal{F}$ induces an
isomorphism}
$$
\cal{F}'_{S'} \to \cal{F}_S \qquad (\text{where $S=\lambda(S')$}).
$$

\begin{proposition}
\label{VI.4.1} Let $f\colon \cal{F} \to \cal{G}$ be an
$\cal{E}$-functor. If~$f$ is fully faithful, then for every
change of base $\cal{E}' \to \cal{E}$, the corresponding functor
$f' \colon \cal{F}' = \cal{F} \times_{\cal{E}} \cal{E}' \to\allowbreak \cal{G}' =
\cal{G} \times_{\cal{E}} \cal{E}'$ is fully faithful.
\end{proposition}

The verification is immediate; more generally, one
can show that every projective limit of
fully faithful functors
(here,~$f$ and the identity functors
in~$\cal{E},\cal{E}'$) is a fully faithful functor.

One notes that the assertion analogous to 4.1,
where
``fully faithful'' is replaced by ``equivalence
of categories'', is false, already for~$\cal{G}=\cal{E}$.
However:

\begin{proposition}
\label{VI.4.2} Let $f \colon \cal{F} \to \cal{G}$ be an
$\cal{E}$-functor. The following conditions are equivalent:
\begin{enumerate}
\item[(i)] There exists an $\cal{E}$-functor $g \colon \cal{G} \to
\cal{F}$ and $\cal{E}$-isomorphisms
$$
gf \isomto \id_{\cal{F}}, \quad fg \isomto \id_{\cal{G}}.
$$
\item[(ii)] For every category $\cal{E}'$ over~$\cal{E}$, the
functor
$$
f' = f \times_{\cal{E}} \cal{E}' \colon \cal{F}' = \cal{F}
\times_{\cal{E}} \cal{E}' \to \cal{G}'=\cal{G} \times_{\cal{E}}
\cal{E}'
$$
is an equivalence of categories.
\item[(iii)]
% original p. 160
$f$ is an equivalence of categories, and for every $S \in
\Ob(\cal{E})$, the functor $f_S \colon \cal{F}_S \to \cal{G}_S$ induced
by~$f$ is an equivalence of categories.
\item[(iii~bis)] $f$ is fully faithful, and for every $S \in
\Ob(\cal{E})$ and every
$\eta\in\Ob(\cal{G}_S)$,
there exists a $\xi \in \Ob(\cal{F}_S)$ and an $S$-isomorphism
$u \colon f(\xi) \to \eta$.
\end{enumerate}
\end{proposition}

\subsubsection*{Proof} Obviously (i) implies that~$f$ is an
equivalence of categories (a notion which is defined by the
same condition, but where the isomorphisms of functors are
not required to be $\cal{E}$-morphisms). On the other hand,
it follows from the functorialities of the preceding no.
that condition (i) is preserved after change of
base~$\cal{E}'\to\cal{E}$. It follows that (i)$\To$(ii).
Obviously (ii)$\To$(iii), for it suffices to take $\cal{E}'
= \cal{E}$ and~$\cal{E}'=\{S\}$. It is even more trivial that
(iii)$\To$(iii~bis), it remains to prove that (iii~bis)$\To$(i). For this, let us choose for every $\eta \in
\Ob(\cal{G})$ a $g(\eta) \in \Ob(\cal{F})$ and an isomorphism
$u(\eta)\colon f(g(\eta))\to\eta$ which is such that~$q(u(\eta)) =
\id_S$, where~$S=q(\eta)$. This is possible thanks to the
second condition (iii~bis). As is well-known and immediate,
the fact that~$f$ is fully faithful implies that~$g$ can in a
unique way be regarded as a functor
from~$\cal{G}$ to~$\cal{F}$, in such a way that the~$u(\eta)$
define a functorial homomorphism (hence an isomorphism)
$u\colon fg \isomto \id_{\cal{G}}$. Moreover, by construction~$g$ is
an $\cal{E}$-functor and~$u$ an $\cal{E}$-homomorphism. To the
preceding data there then corresponds a functorial isomorphism
$v\colon gf \to \id_{\cal{F}}$, defined by the condition
that $f*v=u*f$, and one sees at once that it is likewise an
$\cal{E}$-morphism, qed.

\begin{definition}
\label{VI.4.3} If the preceding conditions are
satisfied, one says that~$f$ is an equivalence of
categories over~$\cal{E}$, or an $\cal{E}$-equivalence.
\end{definition}

\begin{corollary}
\label{VI.4.4} Suppose that the projection functor $p\colon
\cal{F}\to\cal{E}$ is a transportable functor, \ie that for every
isomorphism $\alpha\colon T \to S$ in~$\cal{E}$ and every object~$\xi$
in~$\cal{F}_T$, there exists an isomorphism~$u$ in~$\cal{F}$ with
source~$\xi$ such that~$p(u)=\alpha$. Then every $\cal{E}$-functor
$f\colon \cal{F}\to\cal{G}$ which is an equivalence of
categories, is an $\cal{E}$-equivalence.
\end{corollary}

Follows
% original p. 161
from the criterion (iii~bis).

\begin{corollary}
\label{VI.4.5} Let $f\colon \cal{F}\to\cal{G}$ be an
$\cal{E}$-equivalence. Then for every category~$\cal{H}$
over~$\cal{E}$, the corresponding functors:
\begin{align*}
\SheafHom_{\cal{E}/-}(\cal{G},\cal{H}) & \to
\SheafHom_{\cal{E}/-}(\cal{F},\cal{H}) \\
\SheafHom_{\cal{E}/-}(\cal{H},\cal{F})
& \to \SheafHom_{\cal{E}/-}(\cal{H},\cal{G})
\end{align*}
(\cf no.~\ref{VI.2}) are equivalences of categories.
\end{corollary}

This follows from the criterion (i) by the usual reasoning.

\section{Cartesian morphisms, inverse images, cartesian functors}
\label{VI.5}

Let~$\cal{F}$ be a category over~$\cal{E}$, with
projection functor~$p$.

\begin{definition}
\label{VI.5.1} Consider a morphism
$$
\alpha\colon \eta\to\xi
$$
in~$\cal{F}$, and let
$$
S = p(\xi), \quad T = p(\eta), \quad f=p(\alpha).
$$
One says that~$\alpha$ is a \emph{cartesian morphism}
if for every $\eta' \in \Ob(\cal{F}_T)$ and every $f$-morphism $u\colon
\eta'\to\xi$, there exists a unique $T$-morphism $\overline{u}\colon
\eta'\to\eta$ such that~$u=\alpha\circ\overline{u}$.
\end{definition}

This therefore means that for every $\eta'\in\Ob(\cal{F}_T)$,
the map $v \mto \alpha\circ v$:
\begin{equation}
\tag{i} \Hom_T(\eta',\eta) \to \Hom_f(\eta',\xi)
\end{equation}
is bijective. This also means that the pair $(\eta,\alpha)$
\emph{represents the functor in~$\eta'$} $\cal{F}^{\circ}_T \to
\Ens$ of the right-hand side. If for a given morphism $f\colon T \to S$
in~$\cal{E}$, and a given $\xi \in \Ob(\cal{F}_S)$, there
exists such a pair $(\eta,\alpha)$, \ie a
% original p. 162
cartesian morphism~$\alpha$ in~$\cal{F}$ with target~$\xi$, such that
$p(\alpha)=f$, then~$\eta$ is determined in $\cal{F}_T$
up to unique isomorphism. One then says that \emph{the inverse
image of~$\xi$ by~$f$ exists}, and an object~$\eta$ of~$\cal{F}_T$
equipped with a cartesian $f$-morphism $\alpha\colon \eta \to \xi$ is
called \emph{an inverse image of~$\xi$ by~$f$}.
Often one assumes
such an inverse image chosen whenever it exists ($\cal{F}$
being fixed); one then denotes the inverse image by symbols
such as~$f^*_{\cal{F}}(\xi)$, or simply~$f^*(\xi)$
or~$\xi \times_S T$
when these notations do not lead to
confusion;
the canonical morphism $\alpha\colon \eta \to \xi$ will then be
denoted, in what follows, by~$\alpha_f(\xi)$.
If for every $\xi \in \Ob(\cal{F}_S)$, the inverse image of~$\xi$
by~$f$ exists, one will also say that \emph{the inverse image functor
by~$f$ in~$\cal{F}$ exists}, and~$f^*(\xi)$ then becomes a
\emph{functor covariant in~$\xi$}, from~$\cal{F}_S$ to~$\cal{F}_T$.
This comes from the fact that the right-hand side in (i) depends
covariantly on~$\xi$, \ie more precisely
designates a functor from~$\cal{F}^{\circ}_T \times \cal{F}_S$
to~$\Ens$. This functorial dependence for~$f^*(\xi)$
is made explicit as follows: consider cartesian $f$-morphisms
$$
\alpha\colon \eta \to \xi, \quad \alpha'\colon \eta' \to \xi'
$$
and an $S$-morphism $\lambda\colon \xi \to \xi'$, then there exists a
unique $T$-morphism $\mu\colon \eta\to\eta'$ such that one has
$$
\alpha' \mu = \lambda \alpha
$$
(as follows from the fact that~$\alpha'$ is cartesian).

Let us also note the following immediate fact: consider a
commutative diagram
$$
\xymatrix@C=1.2cm{
\xi\ar[d]_-{\lambda}& \ar[l]_-{\alpha}\eta \ar[d]^-{\mu}\\
\xi' &\ar[l]_-{\alpha'}\eta'
}
$$
in~$\cal{F}$,
% original p. 163
where~$\alpha$ and~$\alpha'$ are $f$-morphisms, and~$\lambda$ an
$S$-isomorphism,~$\mu$ a $T$-isomorphism. \emph{For~$\alpha$
to be cartesian, it is necessary and sufficient that~$\alpha'$
be so}.

\begin{definition}
\label{VI.5.2} An $\cal{E}$-functor $F\colon \cal{F}\to\cal{G}$ is
called a \emph{cartesian functor}
if it transforms cartesian morphisms into cartesian
morphisms. One denotes by $\SheafHom_\cart(\cal{F},\cal{G})$
the full subcategory of
$\SheafHom_{\cal{E}/-}(\cal{F},\cal{G})$ formed of the
cartesian functors.

For example, considering~$\cal{E}$ as a category
over~$\cal{E}$ by means of the identity functor, every morphism
of~$\cal{E}$ is cartesian, hence a cartesian functor
from~$\cal{E}$ to~$\cal{F}$ is a section functor $F\colon
\cal{E}\to\cal{F}$ which transforms every morphism of~$\cal{E}$ into a
cartesian morphism; such a functor is called a \emph{cartesian
section} of~$\cal{F}$ over~$\cal{E}$.
\end{definition}

\begin{proposition}
\label{VI.5.3} \textup{(i)}~A functor $F\colon \cal{F} \to \cal{G}$ which is an
$\cal{E}$-equivalence, is a cartesian functor.
\textup{(ii)}~Let~$F,G$ be two \emph{isomorphic} $\cal{E}$-functors~$\cal{F}\to\cal{G}$. If one is cartesian,
so is the other. \textup{(iii)}~The composite of two cartesian functors
$\cal{F}\to\cal{G}$ and $\cal{G}\to\cal{H}$ is a
cartesian functor.
\end{proposition}

Assertion (iii) is trivial from the definition, (ii)
follows from the remark preceding~\ref{VI.5.2}, (i) follows
easily from the definition and from the criterion~\ref{VI.4.2}~(iii); more
precisely, a morphism~$\alpha$ in~$\cal{F}$ is
cartesian if and only if~$F(\alpha)$ is so.

\begin{corollary}
\label{VI.5.4}
Let $F \colon \cal{F} \to \cal{G}$ be an
$\cal{E}$-equivalence. Then for every category~$\cal{H}$
over~$\cal{E}$, the corresponding functors $G \mto G \circ F$ and
$G \mto F \circ G$ induce equivalences of
categories:
\begin{eqnarray*}
\SheafHom_\cart(\cal{G},\cal{H})\lto{\approx}
\SheafHom_\cart(\cal{F},\cal{H}) \\
\SheafHom_\cart(\cal{H},\cal{F})\lto{\approx}
\SheafHom_\cart(\cal{H}, \cal{G})
\end{eqnarray*}
\end{corollary}

This is deduced in the usual way from~\ref{VI.4.2}
criterion~(i) and from~\ref{VI.5.3} (i) (ii)~(iii).
% original p. 164
One can make more precise that \emph{the $\cal{E}$-functor
$G\colon \cal{G} \to \cal{H}$ is cartesian if and only if $G
\circ F$ is so}, and likewise \emph{an $\cal{E}$-functor $G\colon
\cal{H} \to \cal{F}$ is cartesian if and only if $F \circ G$
is so}.

It follows from~\ref{VI.5.4}~(iii) that if one considers the
subcategory~$\Cat^\cart_{/\cal{E}}$
of~$\Cat_{/\cal{E}}$ whose objects are the same as those of
$\Cat_{/\cal{E}}$, and whose morphisms are the
\emph{cartesian} functors, then one has as in~no.~\ref{VI.2} pairings:
$$
\SheafHom_\cart(\cal{F},\cal{G})\times
\SheafHom_\cart(\cal{G},\cal{H}) \to \SheafHom_\cart(\cal{F},\cal{H})
$$
induced
by those of~no.~\ref{VI.2}, allowing one to consider
$\SheafHom_\cart(\cal{F},\cal{G})$ as a functor in $\cal{F}$,
$\cal{G}$, from the category $\left(\Cat^\cart_{/\cal{E}}
\right)^\circ \times \Cat^\cart_{/\cal{E}}$
to~$\Cat$.
We will need this remark especially for the case where
$\cal{F} = \cal{G}$:

\begin{definition}
\label{VI.5.5}
Let $\cal{F}$ be a category over~$\cal{E}$. One denotes by
$$
\varprojLim\cal{F}/\cal{E}
$$
the category of cartesian $\cal{E}$-functors
$\cal{E}\to\cal{F}$, \ie of cartesian sections of~$\cal{F}$
over~$\cal{E}$.
\end{definition}

According to what has just been said,
$\varprojLim\cal{F}/\cal{E}$ is a functor
in~$\cal{F}$, from the category $\Cat^\cart_{/\cal{E}}$ to the
category $\Cat$.

We will see below the relations between this operation
$\varprojLim$ and the notion of projective limit of
categories, as well as numerous examples.

\section[Fibered categories and prefibered categories]{Fibered categories and prefibered categories. Products and change of base therein}
\label{VI.6}

\begin{definition}
\label{VI.6.1}
A category $\cal{F}$ over~$\cal{E}$ is called a
\emph{fibered category}
(and one then says that the functor $\cal{F}\to\cal{E}$ is
\emph{fibrant})
if it satisfies the following two axioms:
\begin{itemize}
\item{$\Fib_{\mathrm{I}}$}
% original p. 165
For every morphism $f\colon T\to S$ in~$\cal{E}$, the inverse image functor
by $f$ in $\cal{F}$ exists.
\item{$\Fib_{\mathrm{II}}$} The composite of two cartesian morphisms
is cartesian.
\end{itemize}
A category $\cal{F}$ over~$\cal{E}$ satisfying the condition
$\Fib_{\mathrm{I}}$ is called a \emph{prefibered category
over~$\cal{E}$}.
\end{definition}

If~$\cal{F}$ is a fibered category
(\resp prefibered) over~$\cal{E}$, a subcategory
$\cal{G}$ of~$\cal{F}$ is called a \emph{fibered subcategory}
(\resp a \emph{prefibered subcategory}) if it is a
fibered category (\resp prefibered) over~$\cal{E}$, and
if moreover the inclusion functor is cartesian. If for example
$\cal{G}$ is a \emph{full} subcategory of~$\cal{F}$, one
sees that this means that for every morphism $f\colon T\to S$ in
$\cal{E}$ and for every $\xi \in \Ob(\cal{G}_S)$, $f^*_{\cal{F}}(\xi)$
is $T$-isomorphic to an object of~$\cal{G}_T$. Another
interesting case is the following: $\cal{F}$ being a
fibered category over~$\cal{E}$, consider the subcategory
$\cal{G}$ of~$\cal{F}$ having the same objects, and whose morphisms
are the \emph{cartesian} morphisms of~$\cal{F}$; in particular
the morphisms of~$\cal{G}_S$ are the isomorphisms of
$\cal{F}_S$. One sees at once that this is indeed a fibered subcategory
of~$\cal{F}$, for in the bijection
$$
\Hom_T(\eta',\eta) \isomto \Hom_f(\eta',\xi)
$$
relative to a cartesian $f$-morphism $\alpha$ in $\cal{F}$,
to the $T$-isomorphisms of the left-hand side correspond the cartesian
morphisms of the right-hand side. By definition, the cartesian
sections $\cal{E}\to\cal{F}$
then correspond bijectively to arbitrary $\cal{E}$-functors
$\cal{E}\to\cal{G}$ (but one notes that the natural functor
$$
\SheafHom_{\cal{E}/-}(\cal{E},\cal{G})\to
\SheafHom_\cart(\cal{E},\cal{F})=
\varprojLim(\cal{F}/\cal{E})
$$
is faithful, but in general is not fully
faithful, \ie is not an isomorphism).

\begin{remarks}
Let $\cal{F}$ be a category over~$\cal{E}$. The following
conditions are equivalent:~(i)~All morphisms of~$\cal{F}$
are cartesian (ii)~$\cal{F}$~is a fibered category
over~$\cal{E}$, and the $\cal{F}_S$ are groupoids
(\ie every morphism in $\cal{F}_S$ is an isomorphism). One
then says that $\cal{F}$ is a category \emph{fibered in
groupoids}
% original p. 166
over~$\cal{E}$. These are the ones one encounters especially in ``theory
of modules''. If $\cal{E}$ is a groupoid, one shows that
conditions (i) and (ii) are also equivalent to the following:
(iii)~$\cal{F}$ is a groupoid, and the projection functor
$p\colon \cal{F}\to\cal{E}$ is transportable (\cf \ref{VI.4.4}). For
example, if $\cal{E}$ and $\cal{F}$ are groupoids such that
$\Ob(\cal{E})$ and $\Ob(\cal{F})$
are reduced to a point, so that $\cal{E}$ and $\cal{F}$
are defined, up to isomorphism, by groups $E$ and
$F$, and the functor $p\colon \cal{F}\to\cal{E}$ is defined by a
group homomorphism $p\colon F\to E$, then $\cal{F}$ is
fibered over~$\cal{E}$ if and only if $p$ is surjective, \ie if
$p$ defines an \emph{extension} of the group $E$ by the group $G =
\Ker p$.
\end{remarks}

\begin{proposition}
\label{VI.6.2}
Let $F \colon \cal{F} \to \cal{G}$ be an
$\cal{E}$-equivalence. For $\cal{F}$ to be a
fibered category (\resp prefibered) over~$\cal{E}$, it is necessary and
sufficient that $\cal{G}$ be so.
\end{proposition}

Follows easily from the definitions and from the remark
noted above that a morphism
$\alpha$
in $\cal{F}$ is cartesian if and only if $F(\alpha)$ is so.

\begin{proposition}
\label{VI.6.3}
Let $\cal{F}_1$, $\cal{F}_2$ be two categories over~$\cal{E}$, and
let $\alpha = (\alpha_1,\alpha_2)$ be a morphism in
$\cal{F}=\cal{F}_1 \times_{\cal{E}} \cal{F}_2$. For $\alpha$ to be
cartesian, it is necessary and sufficient that its components be so.
\end{proposition}

Indeed, let $\xi_i$ be the target and $\eta_i$ the source of~$\alpha_i$, and
let $f\colon T\to S$ be the morphism of~$\cal{E}$ such that $\alpha_1$ and
$\alpha_2$ are $f$-morphisms. For every
$\eta'=(\eta'_1,\eta'_2)$ in $\cal{F}_T$, one has a commutative
diagram
$$
\xymatrix{ \Hom_T(\eta',\eta) \ar[r] \ar[d] & \Hom_f(\eta',\xi) \ar[d]
\\ \Hom_T(\eta'_1,\eta_1)\times\Hom_T(\eta'_2,\eta_2) \ar[r] &
\Hom_f(\eta'_1,\xi_1)\times \Hom_f(\eta'_2,\xi_2) }
$$
where the vertical arrows are bijections. Thus if one
of the horizontal arrows is a bijection, so is
the other. This already shows that if $\alpha_1$, $\alpha_2$
are cartesian (hence the second horizontal arrow
bijective) then $\alpha$ is so. The converse is seen by setting
in the diagram above $\eta'_i=\eta_i$ whence
$\Hom_T(\eta'_i,\eta_i) \neq \emptyset$, first for $i=2$ which
proves that $\alpha_1$ is cartesian, then for $i=1$ which proves
that $\alpha_2$ is cartesian.

\begin{corollary}
\label{VI.6.4}
Let
% original p. 167
$\cal{F} = \cal{F}_1\times_{\cal{E}}\cal{F}_2$, and let
$F=(F_1,F_2)$ be an $\cal{E}$-functor $\cal{G}\to\cal{F}$. For
$F$
to be cartesian, it is necessary and sufficient that $F_1$ and $F_2$
be so. One thus obtains an isomorphism of categories
$$
\SheafHom_\cart(\cal{G}, \cal{F}_1 \times_{\cal{E}} \cal{F}_2) \isomto
\SheafHom_\cart(\cal{G},\cal{F}_1)
\times\SheafHom_\cart(\cal{G},\cal{F}_2)
$$
and in particular (setting $\cal{G}=\cal{E}$) an isomorphism of
categories
$$
\varprojLim(\cal{F}_1 \times_{\cal{E}}
\cal{F}_2/\cal{E}) \isomto
\varprojLim(\cal{F}_1/\cal{E}) \times
\varprojLim(\cal{F}_2/\cal{E})
$$
\end{corollary}

\begin{corollary}
\label{VI.6.5}
Let $\cal{F}_1$ and $\cal{F}_2$ be two fibered categories
(\resp prefibered) over~$\cal{E}$, then their fibered
product $\cal{F}=\cal{F}_1 \times_{\cal{E}}\cal{F}_2$ is a
fibered category (\resp prefibered) over~$\cal{E}$.
\end{corollary}

These results moreover extend to the case of the fibered
product of an arbitrary family of categories over~$\cal{E}$.

\begin{proposition}
\label{VI.6.6}
Let $\cal{F}$ be a category over~$\cal{E}$, with
projection functor $p$, and let $\lambda \colon \cal{E}' \to \cal{E}$
be a functor, consider
$\cal{F}'=\cal{F}\times_{\cal{E}}\cal{E}'$ as a category over~$\cal{E}'$ by the projection functor $p'=p \times_{\cal{E}}
\id_{\cal{E}'}$. Let $\alpha'$ be a morphism of~$\cal{F}'$; for
$\alpha'$ to be a cartesian morphism, it is necessary and sufficient that
its image $\alpha$ in $\cal{F}$ be so.
\end{proposition}

The proof is immediate and left to the reader.

\begin{corollary}
\label{VI.6.7}
For every cartesian functor $F\colon \cal{F} \to \cal{G}$ of
categories over~$\cal{E}$, the functor
$F'=F\times_{\cal{E}}\cal{E}'$ from
$\cal{F}'=\cal{F}\times_{\cal{E}}\cal{E}'$ to
$\cal{G}'=\cal{G}\times_{\cal{E}}\cal{E}'$ is cartesian.
\end{corollary}

Consequently, the functor $\SheafHom_{\cal{E}}(\cal{F},\cal{G}) \to
\SheafHom_{\cal{E}'}(\cal{F}',\cal{G}')$ considered in
no.~\ref{VI.3} induces a functor
$$
\SheafHom_\cart(\cal{F},\cal{G}) \to
\SheafHom_{\cart}(\cal{F}',\cal{G}')\,;
$$
in other terms, for $\cal{F}$, $\cal{G}$ fixed, \emph{one can
consider}
$$
\SheafHom_\cart(\cal{F} \times_{\cal{E}} \cal{E}',
\cal{G}\times_{\cal{E}} \cal{E}')
$$
\emph{as
% original p. 168
a functor in~$\cal{E}'$, from the category
${\Cat_{/\cal{E}}}^\circ$ to $\Cat$}. If one likewise lets
$\cal{F}$, $\cal{G}$ vary, one finds
a functor from the category ${\Cat_{/\cal{E}}}^\circ \times \big(
\Cat^\cart_{/\cal{E}} \big)^\circ \times \Cat^\cart_{/\cal{E}}$
to~$\Cat$. When one takes into account the isomorphism
$$
\Hom_{\cal{E}'}(\cal{F}',\cal{G}') \isomto
\Hom_\cal{E}(\cal{F}\times_\cal{E}\cal{E}',\cal{G})
$$
envisaged in no.~\ref{VI.3}, then the cartesian $\cal{E}'$-functors
from~$\cal{F}'$ to~$\cal{G}'$ correspond to the $\cal{E}$-functors
$\cal{F} \times_\cal{E} \cal{E}' \to \cal{G}$ which transform every
morphism whose first projection is a cartesian
morphism of~$\cal{F}$, into a cartesian morphism of
$\cal{G}$. Setting $\cal{F}=\cal{E}$, one finds (after change
of notation):

\begin{corollary}
\label{VI.6.8}
$\varprojLim(\cal{F}'/\cal{E}')$ is isomorphic to
the full subcategory of
$\SheafHom_{\cal{E}/-}(\cal{E}',\cal{F})$ formed of the
$\cal{E}$-functors $\cal{E}'\to\cal{F}$ which transform arbitrary
morphisms into cartesian morphisms. In particular, if $\cal{F}$
is a fibered category and if $\tilde{\cal{F}}$ is the
subcategory of~$\cal{F}$ whose morphisms are the
cartesian morphisms of~$\cal{F}$, then one has a bijection
$$
\Ob \varprojLim(\cal{F}'/\cal{E}') \isomto
\Hom_{\cal{E}/-}(\cal{E}',\tilde{\cal{F}})\,.
$$
\end{corollary}

This makes more precise the way in which the expression
$\varprojLim(\cal{F} \times_{\cal{E}} \cal{E}' /
\cal{E}')$ is to be considered as a functor in
$\cal{E}'$ and in $\cal{F}$, from the category
${\Cat_{/\cal{E}}}^\circ \times \Cat^\cart_{/\cal{E}}$ to the
category $\Cat$. One will see later a more complete
functorial dependence with respect to $\cal{E}'$, when
$\cal{F}$ is required to be a fibered category.

\begin{corollary}
\label{VI.6.9}
Let $\cal{F}$ be a fibered category (\resp prefibered)
over~$\cal{E}$, then $\cal{F}'=\cal{F}\times_{\cal{E}}\cal{E}'$ is
a fibered category (\resp prefibered) over~$\cal{E}'$.
\end{corollary}

\begin{proposition}
\label{VI.6.10} Let $\cal{F}$, $\cal{G}$ be
prefibered categories over~$\cal{E}$, $F$ a cartesian $\cal{E}$-functor
from~$\cal{F}$ to $\cal{G}$. For $F$ to be
faithful,
% original p. 169
(\resp fully faithful, \resp an $\cal{E}$-equivalence)
it is necessary and sufficient that for every
$S\in \Ob(\cal{E})$,
the induced functor $F_S\colon \cal{F}_S\to \cal{G}_S$ be
faithful (\resp fully faithful, \resp an equivalence).
\end{proposition}

Immediate proof from the definitions.

To finish this no., we give some properties of
fibered categories, using the axiom $\Fib_{\mathrm{II}}$.
\begin{proposition}
\label{VI.6.11} Let $\cal{F}$ be a prefibered category
over~$\cal{E}$. For $\cal{F}$ to be fibered, it is necessary and
sufficient that it satisfy the following condition:

$\Fib_{\mathrm{II}}'$: Let $\alpha\colon\eta\to\xi$ be a
cartesian morphism in $\cal{F}$ over the morphism $f\colon T\to S$
of~$\cal{E}$. For every morphism $g\colon U\to T$ in $\cal{E}$, and
every
$\zeta\in\Ob(\cal{F}_U)$,
the map $u\mto\alpha\circ u$:
$$
\Hom_g(\zeta,\eta)\to\Hom_{fg}(\zeta,\xi)
$$
is bijective.
\end{proposition}

In other terms, in a \emph{fibered} category
over~$\cal{E}$, cartesian diagrams are characterized
by a property, a priori stronger than that of the
definition (which one obtains by setting $g=\id_T$ in
the preceding statement).
\begin{corollary}
\label{VI.6.12} Let $\cal{F}$ be a category over~$\cal{E}$,
$\alpha$ a morphism in $\cal{F}$. For $\alpha$ to be an
isomorphism, it is necessary that $p(\alpha)=f$ be an isomorphism and that
$\alpha$ be cartesian; the converse is true if $\cal{F}$
is fibered over~$\cal{E}$.
\end{corollary}

Indeed, if $\alpha$ is an isomorphism then so is evidently
$f=p(\alpha)$; for every
$\eta'\in\Ob(\cal{F}_T)$,
the map $u\mto \alpha\circ u$
$$
\Hom(\eta',\eta)\to \Hom(\eta',\xi)
$$
is bijective; as $f$ is an isomorphism, one sees at once
that an
element of the left-hand side is a $T$-morphism if and only
if its image in the right-hand side is an $f$-morphism, hence one thus obtains a bijection
$$
\Hom_T(\eta',\eta)\to\Hom_f(\eta',\xi)
$$
which
% original p. 170
proves the first assertion. Conversely, suppose
that $f$ is an isomorphism and that $\alpha$ satisfies the condition
stated in $\Fib_{\mathrm{II}'}$ (which therefore means,
when $\cal{F}$ is
fibered
over~$\cal{E}$, that $\alpha$ is cartesian), then one sees at
once that for every $\zeta\in\Ob\cal{F}$, the map $u\mto
\alpha\circ u$ from~$\Hom(\zeta,\eta)$ to $\Hom(\zeta,\xi)$ is
bijective, hence $\alpha$ is an isomorphism.
\begin{corollary}
\label{VI.6.13} Let $\alpha\colon\eta\to\xi$ and
$\beta\colon\zeta\to \eta$ be two composable morphisms in the
category $\cal{F}$ fibered over~$\cal{E}$. If $\alpha$ is
cartesian then $\beta$ is so if and only if $\alpha\beta$
is so.
\end{corollary}

One uses the definition of cartesian morphisms in the
strengthened form of~\ref{VI.6.11}.

% ---- en-chunk3.tex ----
